\scanpage{58}
\mspage{14}
C) \uline{Cutting} \deleted{\unsure{of|of the}} along an \inserted{open and closed} subset
$\Sigma'\subset\Sigma_0$ of the submanifold formed by the cutting circles --- this corresponds to a subset $A'$ of $A_0$ (\inserted{bound} edges \deleted{\unsure{ord.|sol.}}), but this time it is not a matter of ``erasing\unsure{''|'} i.e of contracting to a point, but of ``cutting'', which at the level of graphs amounts to cutting an edge at its midpoint (i.e. replacing it by two free edges).

The \uline{marking follows}. \deleted{Here again, for an \unsure{at}}
But here one has to distinguish between the case of an anodyne operation \inserted{at $C$} --- that
where \deleted{there is \unsure{reason to mark|indeed}}
\inserted{$C$ is actually marked} i.e. $C\cap R\ne\phi$ i.e. $C\subset\Sigma^+$ --- which more or less (an ambiguity in a finite set, \deleted{)} \unsure{in near}) can be \unsure{reversed|inverted}, with the help of the product of this operation --- and that where $C$ is not marked,
But it seems that in the way we are going to work
\mspage{57}
